\noindent\textbf{Elección del procedimiento.} La hidrostática da $dp/dz=-\rho g_0$. Sustituyendo $\rho=p/(R_gT)$ y $T=T_0-Lz$, se separa $dp/p=-g_0\,dz/[R_g(T_0-Lz)]$. Integrar desde el nivel del mar da $\ln(p/p_0)=[g_0/(R_gL)]\ln(T/T_0)$. Exponenciar recupera la ley barométrica.

\begin{enumerate}

\item Calcular primero la temperatura en la capa con gradiente.

\[\fitmath{T=T0-L\cdot z}\]

\[\fitmath{T=(288.15)-(0.0065)\cdot (11000)\approx 216.65\,\mathrm{K}}\]

\item El exponente debe ser adimensional.

\[\fitmath{\mathrm{aexp}=\frac{g}{R_g\cdot L}}\]

\[\fitmath{\mathrm{aexp}=\frac{(9.80665)}{(287.05287)\cdot (0.0065)}\approx 5.25588\,}\]

\item Aplicar la potencia al cociente de temperaturas absolutas.

\[\fitmath{p=p0\cdot {\left(\frac{T}{T0}\right)}^{\mathrm{aexp}}}\]

\[\fitmath{p=(101325)\cdot {\left(\frac{(216.65)}{(288.15)}\right)}^{(5.255879813)}\approx 22632.04\,\mathrm{Pa}}\]

\end{enumerate}
